\section{Language-Addressable Spline Action Programs}
\label{sec:method}

Let a VLA policy $\pi_\theta$ receive image observations $o_t$, robot state $s_t$, and language $\ell_t$. A conventional chunking policy predicts a fixed sequence $[a_t,\ldots,a_{t+H-1}]$. We replace only this output target and its decoder; the visual encoder, language model, and flow-matching action expert remain unchanged.

\subsection{Shape--Time Factorization}

For an action segment of duration $T$, define the cumulative commanded end-effector path $p_0=0$ and $p_k=\sum_{j<k}a^{\mathrm{pose}}_j$. We represent this path by a clamped, uniform cubic B-spline
\begin{equation}
    p(u)=\sum_{i=0}^{n-1}B_{i,3}(u)c_i, \qquad u=t/T\in[0,1],
    \label{eq:spline}
\end{equation}
with $n=8$ control points. Endpoints are pinned, $c_0=0$ and $c_{n-1}=p_T$; the interior controls minimize $\sum_{k=0}^{T}\|p(k/T)-p_k\|_2^2$. We precompute the resulting least-squares operators for every allowed $T$, reducing target construction and decoding to batched matrix products. A separate spline fits the raw binary gripper command and is thresholded by sign at decode.

One policy output token represents one pose and gripper control point. An eighth supervised channel contains $\log T$, replicated over the eight tokens and averaged at decode. Thus the action expert emits eight structured tokens instead of 50 waypoint tokens. Targets are standardized per token inside the policy; the external action processor is the identity so that decoded deltas retain their physical meaning. We train with SmolVLA's original conditional flow-matching objective \cite{lipman2023flow,shukor2025smolvla}, applied to these transformed targets.

\subsection{Event-Aligned Targets}

Fixed windows often bisect grasps or contacts. We instead terminate a training target at the first event after a minimum of eight steps: (i) a gripper-sign transition, (ii) two consecutive pose speeds below $0.15$ times the within-window median, or (iii) episode end. If none occurs, the segment is capped at 24 steps (16 for 10-Hz CALVIN). The selected endpoint gives both the path-fitting interval and duration label. This construction makes $T$ a supervised time-to-event rather than an arbitrary chunk size.

At inference, the model predicts $\hat T=\exp(\frac{1}{n}\sum_i\widehat{\log T}_i)$, clamped to the training range. Sampling Eq.~\eqref{eq:spline} at $h$ phases $u_j$ gives positional commands $p(u_j)$; differences $p(u_j)-p(u_{j-1})$ recover executable delta actions. Replanning at $\hat T$ therefore aligns policy queries to predicted manipulation events.

\subsection{Executable Language Clauses}

A long instruction is represented as an ordered program $\mathcal{L}=(\ell^1,\ldots,\ell^M)$. The active clause replaces the whole-task caption in the ordinary VLA language input; no planner, tool library, or architecture change is introduced. For the main two-subgoal study, demonstrations are split at the first gripper release following at least 48 consecutive closed commands. The release remains under clause 1 and the next observation/action pair begins clause 2. All images, states, actions, episode indices, and optimization schedules are otherwise identical to the whole-caption dataset.

During execution a scheduler advances the clause pointer and clears actions generated under the preceding clause. The primary head comparison uses the same outcome-independent clock for waypoint and spline policies: the task-specific median causal switch frame from the training demonstrations (139 and 107 steps). This isolates the action-representation interaction. A second deployment condition advances on the policy's predicted close-to-open gripper event; a separate privileged predicate scheduler is used only for the diagnostic order intervention. The duration channel independently controls when the policy replans within each clause.

\subsection{One Curve, Multiple Physical Clocks}

Because $p$ is parameterized by normalized phase, timing changes need not alter its learned geometry. A monotone time map $\phi:[0,1]\!\rightarrow\![0,1]$ produces $p(\phi(u))$: uniform compression changes overall pace, while interval-selective warping protects slow contact regions. Sampling $p$ at $h=\rho\hat T$ rather than $\hat T$ changes control rate by ratio $\rho$ while preserving endpoints and nominal velocity. If any decoded delta exceeds an actuator bound, feasibility stretch increases $h$ until the bound is met. These operations use one checkpoint and require no additional policy queries or neural inference.
